\documentclass[11pt,a4paper]{article}
\usepackage[british]{babel}
\usepackage[a4paper,margin=2.5cm]{geometry}
\usepackage{booktabs,enumitem,xcolor,listings,float}
\usepackage[T1]{fontenc}
\usepackage{xurl}
\usepackage[hidelinks]{hyperref}
\setlength{\parskip}{4pt}
\setlength{\parindent}{0pt}
\lstset{basicstyle=\ttfamily\small, backgroundcolor=\color{gray!8}, frame=single, rulecolor=\color{gray!40},
  breaklines=true, columns=fullflexible, keepspaces=true, showstringspaces=false}

\title{Configuration Manual\\[4pt]\large Physics-Informed HVAC Energy Prediction under Incomplete Sensor Data}
\author{Sathish Kumar Konda (25190890)\\MSc in Artificial Intelligence, National College of Ireland}
\date{December 2026}

\begin{document}
\maketitle

\section{Purpose}
This manual explains how to set up the environment, obtain the data and reproduce every table and figure of the report. All commands are run from the project folder \texttt{satish-thesis/}.

\section{Hardware}
\begin{table}[H]
\centering\small
\begin{tabular}{lll}
\toprule
Item & Used for the report & Minimum \\
\midrule
Processor & Apple silicon, 8+ cores & any 64-bit x86 or ARM CPU \\
Accelerator & Apple MPS (automatic) & none (CPU works); CUDA is used when present \\
Memory & 16 GB & 8 GB \\
Disk & 1 GB (400 MB for optional raw data) & 200 MB \\
Run time, full experiment & about 45 minutes & about 90 minutes on 4 CPU cores \\
\bottomrule
\end{tabular}
\end{table}

\section{Software}
\begin{table}[H]
\centering\small
\begin{tabular}{ll}
\toprule
Package & Version \\
\midrule
Python & 3.13.9 (3.10 or newer works) \\
numpy / pandas / scipy & 2.4.6 / 2.3.3 / 1.18.0 \\
scikit-learn & 1.5.1 \\
xgboost & 3.2.0 \\
torch & 2.11.0 \\
matplotlib & 3.10.6 \\
pytest, jupyter / nbconvert & any recent version \\
LaTeX (optional, for the report) & TeX Live or TinyTeX with \texttt{latexmk} \\
\bottomrule
\end{tabular}
\end{table}

\section{Set-up}
\subsection{macOS or Linux}
\begin{lstlisting}[language=bash]
git clone https://github.com/Tarunchintakunta/Thesis-2.0.git
cd Thesis-2.0/satish-thesis
python3 -m venv .venv
source .venv/bin/activate
pip install -r requirements.txt
\end{lstlisting}

\subsection{Windows (PowerShell)}
\begin{lstlisting}
git clone https://github.com/Tarunchintakunta/Thesis-2.0.git
cd Thesis-2.0\satish-thesis
py -3 -m venv .venv
.venv\Scripts\Activate.ps1
pip install -r requirements.txt
\end{lstlisting}
For an NVIDIA GPU on Windows or Linux, install the CUDA build of PyTorch from \url{https://pytorch.org} before the other requirements.

\subsection{Google Colab}
Open \texttt{notebooks/physics\_informed\_hvac.ipynb} in Colab (File, Open notebook, GitHub) and choose a GPU runtime if available. The first cell clones only this project with a sparse checkout and installs missing packages. No other step is needed.

\section{Data}
The hourly table used by all experiments is bundled as \texttt{data/bldg59\_hourly.csv.gz}, so the steps below are optional. To rebuild it from the raw LBNL Building 59 files (Dryad, doi:10.7941/D1N33Q):
\begin{enumerate}[nosep]
\item Create a Kaggle API token (Kaggle, Settings, Create New Token) and save it as \texttt{\textasciitilde/.kaggle/kaggle.json}. Dryad itself blocks scripted downloads, so a Kaggle mirror of the same files is used.
\item Run the download script and the preparation step:
\end{enumerate}
\begin{lstlisting}[language=bash]
bash scripts/download_data.sh
python src/prepare_data.py
\end{lstlisting}
This writes \texttt{data/processed/bldg59\_hourly.csv}, the compressed copy, \texttt{results/data\_dictionary.csv} and \texttt{results/data\_profile.json}. Raw files stay in \texttt{data/raw/}, which git ignores.

\section{Running the project}
\subsection{Tests}
\begin{lstlisting}[language=bash]
python -m pytest tests -q
\end{lstlisting}
Expected output: \texttt{19 passed}.

\subsection{Smoke run (2--3 minutes)}
\begin{lstlisting}[language=bash]
python src/experiment.py --quick
\end{lstlisting}
It trains every model for two epochs with one seed and writes files ending in \texttt{\_quick}. Use it to check the installation.

\subsection{Full experiment}
\begin{lstlisting}[language=bash]
python src/experiment.py
python src/analysis.py
\end{lstlisting}
\texttt{experiment.py} runs, for each task (\texttt{estimate}, then \texttt{forecast}), the tree family and then the neural family in separate processes. Options:
\begin{itemize}[nosep]
\item \texttt{--task estimate|forecast} runs one task only;
\item \texttt{--family tree|nn} runs one model family only;
\item \texttt{--stage main|imputation} runs only the main grid or only the imputation sensitivity study.
\end{itemize}
Selected hyperparameters are cached in \texttt{results/chosen\_<task>\_<family>.json}; delete these files to re-tune. \texttt{analysis.py} writes tables to \texttt{results/tables/} (CSV and LaTeX) and figures to \texttt{figures/}.

\subsection{Notebook}
\begin{lstlisting}[language=bash]
jupyter notebook notebooks/physics_informed_hvac.ipynb
\end{lstlisting}
With \texttt{RUN\_FULL = False} (default) it runs the tests, the smoke run and the analysis on the stored results (about 5 minutes). With \texttt{RUN\_FULL = True} it re-runs the full experiment first. To execute it without a browser:
\begin{lstlisting}[language=bash]
jupyter nbconvert --to notebook --execute --inplace notebooks/physics_informed_hvac.ipynb
\end{lstlisting}

\subsection{Report}
\begin{lstlisting}[language=bash]
cd report
latexmk -pdf thesis.tex
\end{lstlisting}

\section{Outputs}
\begin{table}[H]
\centering\small
\begin{tabular}{p{6.4cm}p{8.2cm}}
\toprule
File & Content \\
\midrule
\texttt{results/metrics\_<task>\_<family>.csv} & One row per model, seed, pattern and rate: MAE, RMSE, MAPE, $R^2$, NRMSE, negative-prediction share, temperature RMSE and RC residual. \\
\texttt{results/predictions\_<task>\_<family>.npz} & Test predictions for every model, seed and condition, targets and timestamps. \\
\texttt{results/tuning\_<task>\_<family>.csv} & Every tuning trial and its validation RMSE. \\
\texttt{results/physics\_coefs\_<task>\_nn.csv} & Learned RC coefficients per seed. \\
\texttt{results/plausibility\_*.csv}, \texttt{cost\_*.csv} & Outdoor-temperature response test; training and inference time. \\
\texttt{results/imputation\_estimate\_*.csv} & Mean and CART imputation runs. \\
\texttt{results/tables/}, \texttt{figures/} & Everything shown in the report. \\
\bottomrule
\end{tabular}
\end{table}

\section{Troubleshooting}
\begin{itemize}[nosep]
\item \textbf{Segmentation fault or a hang on macOS when XGBoost and PyTorch are used in one Python process.} Both ship an OpenMP runtime. Always run \texttt{src/experiment.py}, which keeps them in separate processes; do not import \texttt{xgboost} in the same session as a PyTorch training loop.
\item \textbf{MPS not used.} Check \texttt{python -c "import torch; print(torch.backends.mps.is\_available())"}; it needs macOS 12.3+ and an arm64 Python.
\item \textbf{Small numeric differences between runs on GPU.} MPS and CUDA kernels are not bit-deterministic; differences are well inside the seed-to-seed spread reported in the tables.
\item \textbf{\texttt{kaggle: command not found}.} Install it with \texttt{pip install kaggle}, or skip the download and use the bundled table.
\end{itemize}

\end{document}
